\section{Introduction}\label{sec:introduction}

Let $m$ denote Lebesgue measure on $\R$. For a set $A\subseteq\R$, a
\emph{nontrivial affine copy} of $A$ is a set
\[
 x+sA=\{x+sa:a\in A\},\qquad x\in\R,\quad s\in\R\setminus\{0\}.
\]
The set $A$ is \emph{measure universal} if every Lebesgue-measurable subset
of $\R$ of positive measure contains such a copy. The Erd\H{o}s similarity
conjecture asserts that no infinite subset of $\R$ is measure universal;
the problem was posed by Erd\H{o}s in
\cite[Problem~4.33.7*]{Erdos1974Problems}.
Every nonempty finite pattern $F$ is measure universal. Indeed, if a measurable
set contains a subset $E_0$ of finite positive measure, translation
continuity in $L^1$ gives
$m(\bigcap_{a\in F}(E_0-sa))\to m(E_0)>0$ as $s\to0$.
A point in this intersection supplies an affine copy of $F$.
The difficulty therefore lies in keeping infinitely many points in the
same positive-measure set. We consider the single sequence
\[
 D=\{2^{-n}:n\in\N,\ n\ge1\}.
\]

\begin{theorem}\label{thm:main}
For every $\eta\in(0,1)$, there is a compact set $E_\eta\subseteq[0,1]$
with $m(E_\eta)>1-\eta$ such that
\[
 \text{for every }x\in\R\text{ and every }s\in\R\setminus\{0\},
 \quad x+sD\nsubseteq E_\eta.
\]
Equivalently, for every such $x,s$ there is an integer $n\ge1$ with
$x+s2^{-n}\notin E_\eta$.
\end{theorem}

Thus the dyadic sequence is not measure universal. The theorem concerns
one infinite pattern and all of its signed affine copies; the conjecture
for arbitrary infinite sets is not addressed. Compactness and the
near-full-measure conclusion follow directly from the open sets used in
the construction.

\subsection{Related results}

Classical results concern slowly decreasing sequences.
Falconer~\cite{Falconer1984} and Eigen~\cite{Eigen1985} proved
nonuniversality for positive sequences $a_n\downarrow0$ with
$a_{n+1}/a_n\to1$; see also the proof in
\cite[Corollary~1]{HumkeLaczkovich1998}.
Humke and Laczkovich give a finite hitting-set characterization and
further relative-gap criteria \cite[Theorems~1--2]{HumkeLaczkovich1998}.
Kolountzakis's finite-gap criterion uses arbitrarily large finite subsets
of a positive sequence: the minimum gaps, divided by their largest points,
must have negative logarithms growing sublinearly in the cardinalities
\cite[Theorem~3]{Kolountzakis1997}. Chleb\'ik gives a translation-invariant
formulation normalized by the diameter
\cite[Theorem~15]{Chlebik2015}. Dyadic subsets do not satisfy either
condition: for $n\ge2$ dyadic points $a_1>\cdots>a_n$, the minimum gap
is at most $a_{n-1}\le2^{-(n-2)}a_1$, whereas the diameter
$a_1-a_n$ is at least $a_1/2$. Either normalized minimum gap is therefore
at most $2^{-(n-3)}$.

Additive structure gives another family of results. Bourgain proved
that a sum of three infinite sets is nonuniversal
\cite{Bourgain1987}, and Kolountzakis treated sums of two copies of
certain sequences, including $D+D$ \cite[Section~3.4]{Kolountzakis1997}.
The distinction between the dyadic sequence and larger patterns is
important. Mora Cu\'ellar, Iosevich, Kulkarni, Rojas Aravena and Yavicoli
prove nonuniversality for the sum or difference of a nonconstant geometric
sequence tending to zero and an arbitrary infinite set
\cite[Theorem~2]{MoraCuellarEtAl2026Twofold}. Their Introduction explicitly
separates the single dyadic sequence from this two-fold sumset result.
Avoiding a larger pattern does not imply avoidance of a subset of that
pattern.

Iosevich, Kulkarni, Mora Cu\'ellar, Rojas Aravena and Yavicoli prove a
nonuniversality criterion for sets supporting a Rajchman probability
measure, that is, a probability measure whose Fourier transform tends to
zero at infinity \cite[Theorem~1.5]{IKMRY2026Rajchman}. A countable set
cannot support such a measure \cite[Proposition~4.1]{IKMRY2026Rajchman},
so that criterion does not apply to $D$. Their more general
uniform-density criterion on the circle
\cite[Proposition~3.2]{IKMRY2026Rajchman} also fails for
$D\cup\{0\}$: for every integer $k\ge1$, dilation by $2^k$ leaves its
image modulo one unchanged,
and that image stays a distance $1/4$ from $3/4$.

Other notions of size and other transformation classes also lead to
different questions. For certain sequences tending to zero, Cruz, Lai and
Pramanik obtain sets of Hausdorff dimension one avoiding affine copies of
the sequence together with its limit point; the avoiding sets in their
construction have Lebesgue measure zero
\cite[preprint, Corollary~1.2 and Proposition~4.4]{CruzLaiPramanik2023}.
In the other direction,
Feng, Lai and Xiong's bi-Lipschitz embedding theorem applies to the dyadic
sequence \cite[preprint, Theorem~1.1]{FengLaiXiong2024}. Such nonlinear embeddings
are compatible with the affine nonuniversality in Theorem~\ref{thm:main}.
A separate earlier preprint of Cruz, Lai and Pramanik claiming the full
similarity conjecture was withdrawn because of a gap in its
Proposition~3.3 \cite{CruzLaiPramanik2020Withdrawn}.

\subsection{The construction}

Our main step is to construct an open, $1$-periodic set of small density
that meets every $x+tD$ with $x\in\R$ and $t\in[1,2]$. The next section
states this intermediate result precisely. Dyadic rescaling and
reflection then produce a small open set meeting every signed affine
copy of $D$; its complement in $[0,1]$ is the desired compact set.
Periodic blocking sets and summable measure budgets also appear in
\cite[Lemma~3.1 and Proposition~3.2]{IKMRY2026Rajchman}; the finite construction below supplies
the required hitting property for $D$.

Random cell constructions, discretization of the dilation parameter at
a fixed center, and integration over centers occur in
Kolountzakis \cite[Section~3.2]{Kolountzakis1997},
Chleb\'ik \cite[Section~5]{Chlebik2015}, and
Kolountzakis--Papageorgiou
\cite[Section~3.1]{KolountzakisPapageorgiou2025}.
Exceptional centers are repaired by an open cover in Kolountzakis
\cite[Sections~2.2 and~3.2]{Kolountzakis1997}; Chleb\'ik also describes
adding the centers themselves when the pattern contains zero
\cite[Section~3, final remark]{Chlebik2015}.
We use the open-neighborhood form, since $0$ is a limit point of $D$
but is not in $D$.
An earlier attempt to extend the random-cell approach to the dyadic
sequence is described by Solymosi \cite[Section~3]{Solymosi2011Sparse};
that report does not claim a completed construction.
The argument below implements these principles with separated index
windows and a finite routing construction. All estimates needed for the
proof are supplied here.

To construct the periodic set, we first assign each point of the line to
one leaf of a finite ordered tree with $M$ children at each internal
node and $d$ levels. Independent fair binary tables, evaluated on
dyadic cells, determine the route. At each node, the first successful
test chooses a child; if all tests fail, the last child is chosen by
default. Each leaf carries a second table
whose entries select points with a small prescribed probability $p$.
This keeps the expected density small.

At a node, the default child is chosen only when all $M-1$ tested bits
vanish. Although this is rare for large $M$, taking the depth $d$
sufficiently large makes a default somewhere along the center's route
likely. At the first such node, all nondefault choices have been rejected
by the center. Carefully chosen dyadic translates preserve the route to
that node and its earlier rejections, but consult fresh entries at their
current-window resolution. Many nondefault children and many translates
then supply many potential hits.

Two features make those tests useful simultaneously for every normalized
scale. First, an interval of dyadic indices is assigned to each edge, in
the order in which the tree is traversed. Gaps between these intervals
preserve the earlier routing decisions for all but a small set of
centers. Second, each interval is long enough to control every cell
resolution in its child subtree. This bounds the number of different
local tests as the scale varies. We expose the tables in stages, proving
that the terminal entries used by these adaptively chosen routes are
distinct before multiplying their failure probabilities. A finite union
bound then controls all scales in $[1,2]$.

The resulting set may still miss copies centered in a small closed set.
An open neighborhood of those centers repairs every such copy, since
$x+t2^{-n}$ converges to $x$. This last step is what changes an estimate
on exceptional centers into a statement valid at every center.

The windows and their spatial stability are constructed in
Section~\ref{sec:windows}. Section~\ref{sec:routing} defines the random
set and proves the independent-trial estimate.
Section~\ref{sec:scales} handles the continuum of normalized scales and
removes the exceptional centers. Section~\ref{sec:global} completes the
signed-scale deduction of Theorem~\ref{thm:main}.
